% !TeX root = ../tg.tex
\chapter{分割-联邦学习中轻量化客户端训练方案}
\section{无反向传播训练的实现}
在无反向传播训练框架中，模型被分解为多个小型计算单元，每个单元都可以通过一种仅前向传播的机制独立更新。在训练过程中，每次迭代时只对其中一个单元或其对应的输出 logits 有选择性地注入噪声。这种设计确保了扰动对整体损失的影响可以明确归因于目标单元，从而无需反向传播即可实现高效的梯度估计。

值得注意的是，准确的梯度估计需要对同一单元在多次迭代中施加不同噪声样本的重复扰动。只有通过对损失响应各种独立扰动的结果进行聚合，才能使估计的梯度收敛到真实的梯度方向。

该训练方法的核心思想在于观察到：对模型组件引入微小扰动可能会导致最终损失的可测量变化。通过分析扰动与损失变化之间的关系，就可以估计出该组件相对于损失的梯度方向。具体来说，给定一个干净输入 $ x $，令 $ L_{\text{clean}} $ 表示使用模型原始输出计算得到的损失。当一个噪声向量 $ \xi $ 被添加到某个选定单元的输出 logits 或内部参数中时，可以得到扰动后的损失 $ L_{\text{noisy}} $。该单元参数相对于损失的梯度可近似表示为：
\begin{equation}
    \hat{\nabla}_{\theta} L \approx \frac{L_{\text{noisy}} - L_{\text{clean}}}{\varepsilon} \cdot \xi,
\end{equation}

其中 $ \theta $ 表示目标单元的参数，$ \varepsilon $ 是噪声的缩放因子，而 $ \xi $ 是注入的扰动。这个估计的梯度记作 $ \hat{\nabla}_{\theta} L $，随后可用于执行参数更新，通常使用随机梯度下降（SGD）或自适应优化器如 Adam。

扰动可以应用于单元的输出激活值，也可以直接作用于其内部参数。这两种方法代表了从局部变化建模全局损失敏感性的两种根本不同的方式。

在第一种情况下，噪声直接施加在单元的输出激活值上。这种方法类似于函数空间中的有限差分逼近，并能够直观估计单元输出的变化如何影响全局损失。第二种策略则将噪声注入单元本身的权重和偏置参数中。此时，批次中的每个样本都对应一个不同的、经过扰动的单元参数版本，有效地在前向传播过程中模拟了一个模型分布。两种方法都能在不依赖反向传播的前提下实现梯度估计，并可无缝集成到标准深度学习流程中。

重要的是，该方法中使用的噪声并非固定不变，而是与模型参数一同进行学习。我们引入一个可学习的参数向量 $\mathbf{p} \in \mathbb{R}^{C}$，其中 $ C $ 表示该单元的输出通道数。$\mathbf{p}$ 的每一个元素对应一个输出通道，表示所应用噪声尺度的对数值。该参数以某一初始值进行初始化，并在训练过程中根据噪声对损失的影响进行更新。实际的噪声标准差则通过对 $\mathbf{p}$ 进行指数变换获得：
\begin{equation}
   \sigma = e^{\mathbf{p}},
\end{equation}

该表达式确保了噪声尺度的正定性，并支持基于梯度的优化。这种自适应的噪声缩放机制使模型能够根据每个单元的敏感性自动调整扰动幅度，从而实现更稳定且更精确的梯度估计。

此外，该方法支持多种噪声采样分布，包括标准高斯噪声以及低差异序列。这些拟蒙特卡洛方法相比纯随机采样能更均匀地覆盖输入空间，从而提高梯度估计的可靠性。

通过消除对反向传播的依赖，无反向传播训练为在内存受限或架构限制导致传统梯度计算不可行的场景下训练深度神经网络提供了一种有前景的替代方案。

\section{无反向传播训练对分割-联邦学习的性能提升}
我们考虑一个包含一个中心服务器和 $ n $ 个边缘客户端的分裂学习系统，其中每个客户端 $ n \in \{1, 2, \dots, K\} $ 拥有一个本地数据集 $ \mathcal{D}_n = \{(\mathbf{x}_{ni}, y_{ni})\}_{i=1}^{|\mathcal{D}_n|} $，其中 $ \mathbf{x}_{ki} \in \mathbb{R}^d $ 是输入特征向量，$ y_{ki} \in \mathbb{R} $ 是对应的标签。一个由参数 $ \boldsymbol{\omega} \in \mathbb{R}^p $ 表示的神经网络模型被划分为两个部分：客户端子模型 $ \boldsymbol{\omega}_c: \mathbb{R}^d \rightarrow \mathbb{R}^m $ 和服务端子模型 $ \boldsymbol{\omega}_s: \mathbb{R}^m \rightarrow \mathbb{R} $，使得对输入 $ \mathbf{x} $ 的整体预测表示为：
\begin{equation}
    \hat{y} = f_{\boldsymbol{\omega}_s}(f_{\boldsymbol{\omega}_c}(\mathbf{x})) = f_{\boldsymbol{\theta}}(\mathbf{x}),
\end{equation}


其中 $ \boldsymbol{\omega} = [\boldsymbol{\omega}_c^\top, \boldsymbol{\omega}_s^\top]^\top $ 表示客户端与服务端参数的拼接。

在基于标准反向传播的训练中，每个客户端通过在一个小批量数据上进行前向和后向传播，计算精确梯度 $ \nabla_{\boldsymbol{\omega}_c} L(f_{\boldsymbol{\omega}}(\mathbf{x}), y) $，其中 $ L(\cdot, \cdot) $ 表示损失函数。该过程需要在前向传播过程中存储计算图，导致内存消耗为 $ M_{bp} $。此过程的计算复杂度记为 $ C_{bp} $，每轮迭代中客户端与服务器之间的通信开销为 $ T_{bp} $。

相比之下，在将无反向传播的梯度估计方法集成到分裂学习中的情况下，客户端 $n$ 仅执行前向传播来计算中间表示：
\begin{equation}
    \mathbf{z}_n = f_{\boldsymbol{\omega}_c}(\mathbf{x}_n),
\end{equation}


然后将其传输至服务器。服务器随后计算最终输出 $ \hat{y}_n = f_{\boldsymbol{\omega}_s}(\mathbf{z}_n) $ 并评估损失值 $ L_n = L(\hat{y}_n, y_n) $。该损失值返回给客户端，客户端随后使用无反向传播方法估计梯度 $ \tilde{\nabla}_{\boldsymbol{\theta}_c} L_n $。由于未在本地执行后向传播，前向计算期间的内存占用减少至大约一半，即：
\begin{equation}
    M_{ulr} \approx 0.5 M_{bp},
\end{equation}


因为无需存储用于梯度计算的计算图。

此外，通信开销也减半，因为只有中间输出 $ \mathbf{z}_n $ 从客户端上传到服务器，而只有一个标量 $ L_n $ 从服务器下载到客户端。因此，
\begin{equation}
    T_{ulr} \approx 0.5 T_{bp}.
\end{equation}

\section{实验结果}
本实验在 CIFAR-10 数据集上，基于 ResNet-5 和 VGG-8 两种模型，评估了在分割式联邦学习框架中应用无反向传播训练方法的效果。实验中共有 10 个客户端参与。如表 \ref{tab:bpfree} 所示，采用无反向传播方法进行训练的分割式联邦学习，在测试准确率上接近于集中式训练下使用相同训练策略的结果。
\begin{table}[t]
\centering
\caption{无反向传播训练方法中分割-联邦学习与集中式学习的准确率比较}
\label{tab:bpfree n=10}
\begin{tabular}{c|cc}
\hline
\textbf{训练框架} & \textbf{模型} & \textbf{准确率 (\%)} \\ \hline
集中式学习 & ResNet-5 & 64.80 \\ \cline{2-3}
& VGG-8 & 74.30 \\ \hline
分割-联邦学习 & ResNet-5 & 62.22 \\ \cline{2-3}
& VGG-8 & 74.53 \\ \hline
\end{tabular}\label{tab:bpfree}
\end{table}

\section{未来研究计划}
尽管本研究已初步验证了无反向传播训练方法在分割-联邦学习框架下的可行性，并在ResNet-5和VGG-8等模型上取得了接近传统基于反向传播训练的准确率表现，但该方向仍处于探索阶段，具有较大的拓展空间。未来的研究将从以下几个方面深入展开：

首先，我们将继续探索无反向传播训练方法在更多类型神经网络模型上的适用性，包括但不限于Transformer、轻量级CNN结构（如MobileNet、SqueezeNet）以及图神经网络（GNN）。同时，将在更广泛的数据集上进行实验评估，涵盖图像分类、语义分割等任务，以验证该方法在不同任务场景下的泛化能力与稳定性。

其次，为了推动该技术在边缘设备上的实际部署，我们计划使用基于树莓派的测试平台对应用无反向传播训练的分割-联邦学习算法进行全面的性能评估。重点分析其在资源受限环境下的能耗、内存占用、通信负载等关键指标，为后续优化提供数据支持。通过对比标准分割-联邦学习训练流程，量化无反向传播机制在降低客户端负担方面的实际效益。

此外，针对当前无反向传播方法在高维参数空间中梯度估计精度较低的问题，我们将进一步优化其噪声注入策略与梯度估计方式。例如引入自适应噪声缩放机制、多扰动梯度融合策略，或结合低方差采样技术（如拟蒙特卡洛方法），以提升梯度估计的稳定性和收敛速度。目标是在不牺牲模型准确率的前提下，显著降低客户端的计算与通信开销。

最终，我们希望构建一个高效、稳健且适用于边缘计算场景的分割-联邦学习系统，在保障隐私安全的同时实现高精度、低功耗、低延迟的协同训练模式，为医疗健康、智能物联网等对资源敏感且数据异构性强的应用领域提供技术支持。

